\subsection{把胜利留在地方：西南的军镇与海峡的往返}

一六六一年，\eventanchor{EQ-1661-KANGXI-ACCESSION}{康熙即位（1661）}时，北京交到八岁玄烨和辅政大臣手中的，不只是宫廷的人事安排，还有仍在西南、海峡与各省军镇之间迟缓传递的军情。北京能够发出任命和处分，仍须经旗务、部院、督抚与将领把命令送进不同战区；接到的回报又会被行军、季节与呈报者的利害改变。次年，\eventanchor{EQ-1662-YONGLI-EXECUTION}{永历帝被处死（1662）}，南明最后一个皇帝名义中心消失，但朱由榔自缅甸方向被押回的路线，已把吴三桂集团、清军、缅甸宫廷与边境通行牵在一起。处决是一个政治节点，不是云南西部的军队、流亡者、贸易与安置问题同时终结的时刻。

康熙十二年春，撤藩的题本从北京发出时，云南、广东和福建早已不是可以在部议中随意腾挪的空格。吴三桂、尚可喜、耿氏家族所维系的，是入关征服后留在当地的军队、粮饷渠道、亲属与旧部的前途，也包括州县官、商人和乡绅不得不同这些军镇打交道的日常。中央希望把长期的驻军财政与任命权收回，首先碰到的却是很具体的问题：哪一处仓粮可供转运，哪一路驿道能让新调来的部队通过，原先领饷的人由谁安置。因而\eventanchor{EQ-1673-FEUDATORY-ABOLITION}{撤藩决定（1673）}不是一道命令从北京落到三处属地的单线故事，而是把这些已被嵌入地方的安排骤然推入不确定之中。

吴三桂在云南举兵后，西南的战争也不能只按朝廷与一位将领的胜负来读。滇、黔的山路、渡口和城镇把军队、运夫、商旅、逃难者与地方守备卷在一起；一支部队能否继续向前，取决于粮草是否赶得上、地方是否还能征到人力，也取决于守城者、降将和州县官在消息尚未明朗时如何选择。清军最终进入昆明，改变了吴氏政权的军事前景，却没有自动消除驻军的开销、垦复的压力和战后重建税源的困难。把战争写成若干城池颜色的转换，会遗漏最持久的一层后果：新政权仍须在被扰乱的道路上运输，在经历征发的人群中重新建立可征可管的关系。

官修的战事汇编和实录擅长保存北京收到的捷报、招降、任免与奖惩，因而能让我们追踪朝廷怎样把战局编排为凯旋的次序。它们不能替代军需题本、府县志和地方契约所偶然留下的运粮、逃徙、复业和讼案，更不能据此计算每一地的损失或所有人的立场。所谓“降附”在军报中是一种行政结论；对一名改换主人的军官、一个要保护田土的家族或一个被军队经过的村落而言，它可能分别意味着领饷、避祸、被征发或暂时等待。材料能使这些差异局部可见，却不足以把西南居民写成整齐的一种反应。

海峡的难题则由船、风和港口把同样的行政愿望拆散。迁界意在截断沿岸对郑氏的接济，改变的却是闽粤渔户出海、盐货转运、耕地利用与亲族往来的条件；复界和海禁的调整也不能仅凭一纸上谕推定为同时发生的生活转变。沿海村落有人迁离，有人在限制松动后设法回返，也有人把生计移往别处。地方志常在后来用恢复的语言回顾这些变化，官府命令则记录其规定而非每一次迁移；两者都应与港口账册、契约和航海记录并读，才可知道哪些地点、哪些群体留下了可说的痕迹。

\subsection{从热兰遮到澎湖：岛内土地、海上补给与接收的限度}

一六六一年，\eventanchor{EQ-1661-ZHENG-TAIWAN-LANDING}{郑成功登陆台湾（1661）}，热兰遮城外的围困把海上争夺压缩为船只何时到港、粮食能否支撑、港湾是否可用的连续计算。郑氏需要一个能供给军队并容纳其政治组织的基地；荷兰东印度公司的据点同时嵌在贸易、税收与台湾西南既有的居民关系里。登陆和围城可以确定为军事行动的开端，却不能由此替汉人移民、不同原住民社群、港口商人或荷兰雇员指定同一立场。荷兰城堡日志、清方海防军报和后出的郑氏叙事各自看见不同的敌人、盟友与损失；兵数、伤亡以及各社如何卷入，仍须在这些不相称的记录之间谨慎停留。

郑成功去世后，大陆据点与台湾之间的取舍变得更急迫。\eventanchor{EQ-1663-XIAMEN-JINMEN}{厦门、金门易手（1663）}，并不等于清廷已经在海面上建立无缝的控制；它使郑氏主力退向台湾，也使闽南的岛屿、港口和沿岸村落面对新的航路与征发压力。熟悉水面的降将、地方船户和能否取得船粮，都比地图上两座岛的归属更接近当时的胜负条件。清方军报和荷兰方面的消息可以互相钉住据点的易手，却难以替沿海居民说明一次禁令如何影响其出海，或替郑氏内部还原每一笔补给从何而来。

这种不确定性在\eventanchor{EQ-1681-ZHENG-JING-DEATH}{郑经去世（1681）}后进入东宁的继承问题。郑克塽取得位置，意味着郑氏的宗族、将领与文官已在岛内重新排列；北京得到的情报能够显示清廷怎样估量这个窗口，却不能把郑氏后出的叙事和清方侦报拼接成一场没有空白的宫廷政变。继承冲突削弱了对手协调船队与人力的能力，但不足以单独解释两年后澎湖的结局。福建水师还必须筹集舰船、等待适航的季节、维持兵粮，并判断沿海据点能否供接应。郑氏的战败和投降改变了海峡两岸的军事关系；清廷在台湾设府，是把福建的行政链向东延伸的起点，而非一枚印信立刻重写岛内土地、渡海限制与不同社群关系的终点。

因此，所谓海疆的“接收”必须拆成几条速度不同的线来看。清廷可以任命官员、部署驻防、规定渡台与贸易；船户和移民仍要面对风候、航程、债务与港口机会，岛内村社则在垦地、水源、税役和与官府交涉的具体处境中作出反应。原住民社群留下的直接自述尤为稀少，殖民官员、郑氏与清朝文书又都带有土地和治理的诉求。它们足以提示行政接收改变了谁可以征税、谁能够通航，却不能让我们替每一处聚落断定接受、抵抗或融入。海峡不是陆地政令的末端，而是财政、人群流动和军队补给不断相互改写的空间。

\subsection{草原、河流与高原：多语关系不是一张完成的疆图}

北方的秩序也以迁徙和供给开始，而不是以会盟仪式结束。噶尔丹进攻喀尔喀之后，部众、牲畜与领主向南移动，清廷所面对的既是军事求援，也是安置、赈济、水草和道路的问题。台站要转运消息与物资，草原上的关系却不能由北京单方面排定：请求保护的领主、原有的牧地使用者、寺院网络和负责调度的人，都在同一片季节性移动的空间里。兵力与迁徙人数在不同战报中常有夸张或省略；把它们化为一个精确规模，反而会掩盖记录者究竟是在报功、请饷还是争取援助。

在这样的背景下，\eventanchor{EQ-1691-DOLONNOR}{多伦诺尔会盟（1691）}把清廷的军事保护、封爵礼仪、盟旗编制和物资调度连接起来。会盟能够确证清廷希望以何种名义组织与喀尔喀诸部的关系；理藩院的满、蒙文档案又能使后续的赈济、编制和事务往来局部可见。但“归附”不是可以套在所有人身上的完成式。牧地怎样使用、部际纠纷由谁居中、领主保留何种选择，仍须沿着不同语言的文书和地方记忆逐项追踪。多伦的仪式若脱离这些问题，只会把一段谈判关系误读为地图上已经涂好的边界。

一六八九年尼布楚的交涉提供了另一种限度。河名、山名和地名被写入满文、拉丁文、俄文等文本，是使双方能够暂时结束阿尔巴津危机的办法；翻译、地物对应和使团行程本身就是谈判的一部分。条约的文字可证明签署者承认了怎样的程序和措辞，不能自行量出一条现代精确线，也不能证明河岸的猎人、商旅或巡查者已经知道并接受同一种分类。清方实录所记的批准与赏赐、俄方使团文书所记的交涉、黑龙江地方志及驿路和贸易记录所见的实际往来，应当并置：前两者说明两个国家怎样叙述一次外交成果，后两者才可能让边界在何时、何地被实践的问题显出轮廓。

通往青海、西藏的关系同样不能由册封、驻军或后来的地图倒推为稳定控制。寺院、地方政教力量、蒙古盟军、商旅与沿途的向导、驮运者共同决定道路在某个季节是否可通；军队进入拉萨前后，粮秣、病损与地方合作都限制着它能实行什么。清廷驻防和理藩院文书留下的是一条行政视线，藏文寺院文书和地方传统则有自己的时间线，且彼此保存并不均衡。把它们并读，不是为了制造一个无缝的“真实全景”，而是为了不把其中任何一种制度语言误作所有行动者的经验。

\subsection{把材料的视野留在它能够到达的地方}

从北京望出去，早期清朝似乎由诏令、任命、盟约和凯旋组成；从军需与地方材料望回去，它又是许多不稳定的接口。部院可以下令撤藩或设府，却要通过地方官、将领、旗务机构和中介人员获得粮、船、马匹与消息。边地的行政不是命令距离京城较远的简化版：在云南，它受山路、仓储和战后安置牵制；在台湾，它受季风、港湾、渡海限制和土地关系牵制；在草原与高原，它受牧地、寺院、翻译和季节性道路牵制。不同限制不是中央权力的注脚，而是决定政策在何处改变形状的条件。

《平定三逆方略》和《清圣祖实录》使人看见朝廷如何排出招降、军功与处分的顺序。题本、地方志、郑氏材料以及满、蒙、藏文档案则让供给、迁徙和地方关系在若干节点露出侧面。任何一类材料都不能独自解决另一个层次的问题：一份上谕证明命令曾被发布，不证明村社如何执行；一份条约证明文字曾被同意，不证明地名和界线已经被共同实践；一条地方记录能够照亮一处道路或港口，也不能替整个区域发言。兵额、战果、人口损失和航运规模尤其需要逐项注明记录者、地域与目的，并与对手、地方和跨境材料互校。

这也改变了我们阅读“平定”“归附”或“安民”这类官府语言的方式。它们首先说明国家试图把人、地和关系纳入何种行政分类，而不是居民、盟旗、寺院或绿洲社会已经毫无保留地接受了这种分类。钱穆关于制度如何把中央意图送入边地的提问，能帮助我们追问这些接口；吕思勉关于社会经济条件的提问，则提醒我们回到河流两岸的贸易、牧猎、移徙与驻防。两种提问都不能代替史料的限度：在直接证词缺失处，叙事应保留不知道谁怎样选择的空间，而不以帝国的文书替他们补写回答。
